\chapter{Спецификация Version IV.SM}

Version IV.SM открывается не для поиска ещё одного красивого числа. Её
задача --- построить одну архитектуру, которая имеет наблюдаемый набор полей,
корректно переносится между масштабами и выдерживает заранее объявленные
физические проверки. Численное совпадение после выбора специальных
коэффициентов не считается замыканием.

\section{Исходный диагноз}

Том III.H дал математически согласованную скрытую
\(U(1)\)-теорию, но её минимальное спектральное завершение точно отделено от
Стандартной модели. Поэтому скрытые отношения
\begin{equation}
 g_h^2=\frac{3}{8},
 \qquad
 m_s:m_A:m_f=2:\sqrt3:1,
 \qquad
 B_0=\frac{67}{64\pi^2}
\end{equation}
нельзя сравнивать с массами и связями известных частиц как предсказания
наблюдаемого мира. Version IV.SM должна изменить representation-level
архитектуру, а не переименовать скрытые моды.

\section{Минимальная карта наблюдаемого мира}

Первый обязательный readout имеет вид
\begin{equation}
 \begin{aligned}
 \mathcal R_{\rm SM}:&
 (\mathcal A_F,H_F,D_F,J_F,\gamma_F)\\
 &\longrightarrow
 \left\{
 \begin{array}{l}
 SU(3)_c\times SU(2)_L\times U(1)_Y,\\
 Q_L,u_R,d_R,L_L,e_R,\nu_R\ \text{или обоснованное отсутствие }\nu_R,\\
 H,\ \text{калибровочные вершины и Yukawa-граф},\\
 \text{три поколения или механизм их появления}
 \end{array}
 \right\}.
 \end{aligned}
\end{equation}
Карта должна следовать из представления конечной алгебры, реальной структуры
и допустимых внутренних флуктуаций. Таблица квантовых чисел может быть
временно заморожена как эмпирическая цель, но тогда она не является выводом
теории.

\section{Что означает ``всё совпало''}

Физическое замыкание разбивается на независимые уровни.

\begin{center}
\small
\begin{longtable}{>{\raggedright\arraybackslash}p{0.22\textwidth}>{\raggedright\arraybackslash}p{0.48\textwidth}>{\raggedright\arraybackslash}p{0.20\textwidth}}
\toprule
Уровень & Обязательный результат & Допустимый статус \\
\midrule
\endhead
Представления & полный SM charge ledger, chirality, anomaly cancellation, Higgs и neutrino graph & строгий вывод или явно маркированный input \\
Действие & один gauge-invariant и BRST-consistent low-energy functional & строгий вывод \\
Нормировки & относительные \(g_1,g_2,g_3\), scalar normalization и Yukawa measure из одной trace convention & строгий вывод \\
Масштаб & не более одного объявленного dimensionful train input & калибровка \\
RG & перенос от matching scale до laboratory scales с thresholds и scheme ledger & воспроизводимый расчёт \\
Электрослабый сектор & \(M_W,M_Z,M_H\), \(G_F\), \(\sin^2\theta_W\) и электромагнитная связь & минимум две blind-величины \\
Сильный сектор & \(\alpha_s\) и согласованный цветовой спектр & blind-величина \\
Фермионы & массы, CKM и CP-инвариант без независимой настройки каждой наблюдаемой & разделённый train/blind тест \\
Нейтрино & две mass splittings, mixing pattern и mass ordering branch & минимум одна blind-величина \\
Мост томов & один functional origin для наблюдаемого действия и допустимой части геометрического ядра & строгий вывод \\
\bottomrule
\end{longtable}
\end{center}

\begin{protocol}[Definition of Done Version IV.SM]
Version IV.SM получает статус физически замкнутой только если одновременно:
\begin{enumerate}
 \item конечная алгебра и её представление фиксированы до численного fit;
 \item все локальные gauge и mixed gravitational anomalies сокращаются;
 \item из одного bosonic/fermionic measure следуют относительные
 нормировки gauge, scalar и Yukawa-блоков;
 \item после не более чем одного scale input выполнен RG-перенос;
 \item не менее двух dimensionless observables, не использованных при
 построении модели, проходят заранее зафиксированные tolerance windows;
 \item хотя бы одна blind-проверка относится к EW/QCD, а другая --- к
 fermion или neutrino сектору;
 \item перечислены все неудачные наблюдаемые и запрещено менять модель после
 раскрытия blind-набора без открытия новой версии;
 \item результат воспроизводится независимым машинным pipeline.
\end{enumerate}
\end{protocol}

\section{Порядок работы}

Работа проводится через kill-gates:
\begin{enumerate}
 \item representation gate: получить наблюдаемый charge ledger;
 \item anomaly gate: исключить несогласованные chiral modules;
 \item spectral-normalization gate: вывести отношения \(g_1:g_2:g_3\);
 \item Higgs--Yukawa gate: получить scalar potential и mass graph;
 \item RG gate: перенести matching relations к низким энергиям;
 \item two-sector blind gate: одновременно проверить gauge и matter sectors;
 \item cross-tome gate: проверить, существует ли общий functional bridge.
\end{enumerate}

\begin{nogocriterion}[Запрет числового спасения]
Если после раскрытия blind-наблюдаемой вводится новый коэффициент, меняется
representation, выбирается новый threshold или переназначается физический
смысл внутренней моды, соответствующая версия считается проваленной. Близкое
совпадение отдельного числа не компенсирует неверный спектр частиц,
аномалию или провал независимого сектора.
\end{nogocriterion}

\begin{proposition}[Необходимое изменение архитектуры]
Exact direct-sum completion Тома III.H не может выполнить Definition of Done
Version IV.SM, поскольку не содержит bi-charged connector states и не
порождает ненулевой nongravitational portal. Следовательно, наблюдаемое
замыкание требует новой конечной алгебры или нового недиагонального
bimodule, а не выбора масштаба \(\chi\).
\end{proposition}